\chapter{Toegepaste groepentheorie}\label{ch:b3:group-theory-applied}

Koolstofdioxide maakt enkele honderden moleculen uit op elk miljoen
moleculen lucht, stikstof en zuurstof vrijwel al de rest; toch is het
koolstofdioxide, en niet stikstof of zuurstof, dat het infraroodlicht
absorbeert dat de warme bodem naar de ruimte uitstraalt. Een molecuul
absorbeert infraroodlicht alleen via een vibratie die zijn dipoolmoment
verandert, en of een vibratie dat doet, wordt uitsluitend door de
symmetrie bepaald: de enige vibratie van stikstof kan het niet, twee van
de vier vibraties van koolstofdioxide wel. Dit hoofdstuk maakt van de
karaktertabellen van
\cref{ch:b3:point-groups} werktuigen: het reduceert representaties,
bouwt symmetrieaangepaste orbitalen, telt en benoemt vibraties, en leidt
de selectieregels van infrarood- en ramanspectroscopie af.

\begin{recall}
\Cref{ch:b3:point-groups} definieerde \omterm{def:b3:point-groups:operation}{symmetrieoperaties}, \omterm{def:b3:point-groups:point-group}{puntgroepen},
klassen, representaties, \omterm{def:b3:point-groups:representation}{karakters} en karaktertabellen. Het deel van
jaar~2 bouwde de moleculaire orbitalen van \ce{H2O}, \ce{NH3} en \ce{CH4}
op uit fragmentorbitalen en symmetrieaangepaste combinaties van de
$1s$-orbitalen van waterstof, en benoemde ze met labels die hier
afgeleid worden. Het deel van jaar~1 las infraroodspectra in golfgetallen
en definieerde de polariseerbaarheid van een molecuul.
\end{recall}

\section{Een representatie reduceren}

Elke verzameling functies of vectoren die aan een molecuul vastzit,
draagt een representatie van zijn \omterm{def:b3:point-groups:point-group}{puntgroep}. De meeste zijn reducibel.

\begin{definition}[Reducibele representatie, totaalsymmetrische representatie]\label{def:b3:group-theory-applied:reducible}
Een \emph{reducibele representatie}\index{reducibele representatie} is een
representatie waarvan alle matrices door één basisverandering in dezelfde
blokdiagonale vorm gebracht kunnen worden; ze is dan een som van
\omterm{def:b3:point-groups:irrep}{irreducibele representaties}, geschreven als $\Gamma = \sum_in_i\Gamma_i$. De
\omterm{def:b3:point-groups:irrep}{irreducibele representatie} waarvan alle \omterm{def:b3:point-groups:representation}{karakters} 1 zijn, is de
\emph{totaalsymmetrische
representatie}\index{totaalsymmetrische representatie} ($A_1$, $A_{1g}$,
$A_1'$\dots).
\end{definition}

\begin{theorem}[Grote orthogonaliteitsstelling]\label{thm:b3:group-theory-applied:great-orthogonality}
Voor twee \omterm{def:b3:point-groups:irrep}{irreducibele representaties} $\Gamma_i$, $\Gamma_j$ met dimensies $d_i$,
$d_j$ van een groep van orde $h$, geschreven met unitaire matrices, geldt
\[
\sum_R D_i(R)_{mn}^*\,D_j(R)_{m'n'} = \frac{h}{d_i}\,\delta_{ij}\,\delta_{mm'}\,\delta_{nn'}.
\]
\end{theorem}
\admitted

Het bewijs, met de lemma's van Schur, hoort thuis in meer gevorderde
cursussen; de chemie gebruikt de gevolgen voor de \omterm{def:b3:point-groups:representation}{karakters}, en elke
\omterm{def:b3:point-groups:irrep}{karaktertabel} in dit boek wordt daaraan getoetst.

\begin{corollary}[Orthogonaliteit van de karakters]\label{cor:b3:group-theory-applied:characters}
Met $g_c$ het aantal operaties in klasse $c$ is
\[
\sum_cg_c\,\chi_i(c)^*\chi_j(c) = h\,\delta_{ij},
\]
en het aantal \omterm{def:b3:point-groups:irrep}{irreducibele representaties} is gelijk aan het aantal
klassen, met $\sum_id_i^2 = h$.
\end{corollary}

\begin{proof}
Stel $m = n$ en $m' = n'$ in de stelling en sommeer over $m$ en $m'$: het
linkerlid wordt $\sum_R\chi_i(R)^*\chi_j(R)$, het rechterlid $\frac{h}{d_i}\delta_{ij}
\sum_{m}\delta_{mm} = h\delta_{ij}$; groepeer de operaties per klasse. De
rijen van de tabel zijn dus orthogonale vectoren in een ruimte waarvan de
dimensie het aantal klassen is, zodat er hoogstens evenveel \omterm{def:b3:point-groups:irrep}{irreducibele
representaties} als klassen zijn; de gelijkheid en $\sum d_i^2 = h$ volgen
uit de toepassing van de stelling op de reguliere representatie
(oefening 10).
\end{proof}

\begin{theorem}[De reductieformule]\label{thm:b3:group-theory-applied:reduction}
Een representatie $\Gamma$ met \omterm{def:b3:point-groups:representation}{karakters} $\chi(c)$ bevat de \omterm{def:b3:point-groups:irrep}{irreducibele
representatie} $\Gamma_i$ precies
\[
n_i = \frac1h\sum_cg_c\,\chi(c)\,\chi_i(c)^*
\]
keer. Dit is de \emph{reductieformule}\index{reductieformule}.
\end{theorem}

\begin{proof}
Omdat $\Gamma = \sum_jn_j\Gamma_j$, zijn haar \omterm{def:b3:point-groups:representation}{karakters} $\chi(c) = \sum_jn_j\chi_j(c)$.
Vermenigvuldig met $g_c\chi_i(c)^*$, sommeer over de klassen en gebruik het
gevolg:
$\sum_cg_c\chi(c)\chi_i(c)^* = \sum_jn_jh\delta_{ij} = hn_i$.
\end{proof}

\begin{example}[De waterstoforbitalen van ammoniak]\label{ex:b3:group-theory-applied:nh3}
Neem de drie $1s$-orbitalen van de waterstofatomen van \ce{NH3} als basis. $E$
laat ze alle drie op hun plaats ($\chi = 3$), $C_3$ verplaatst ze alle drie
($\chi = 0$), een
$\sigma_v$ laat er één op zijn plaats en verwisselt er twee ($\chi = 1$). Met de
tabel van $C_{3v}$:
$n_{A_1} = \frac16(3 + 0 + 3\times1) = 1$, $n_{A_2} = \frac16(3 + 0 - 3) = 0$,
$n_E = \frac16(6 + 0 + 0) = 1$. Dus $\Gamma_{\mathrm H} = A_1 + E$: de drie
waterstofatomen geven één totaalsymmetrische combinatie en een ontaard paar.
\end{example}

\begin{method}[Een representatie reduceren]\label{met:b3:group-theory-applied:reduce}
\begin{enumerate}
  \item Zoek voor elke klasse het \omterm{def:b3:point-groups:representation}{karakter}: voor een verzameling functies
        gecentreerd op atomen is dat het aantal functies dat op zijn plaats
        blijft, elk geteld met het teken dat het krijgt.
  \item Pas de \omterm{thm:b3:group-theory-applied:reduction}{reductieformule} toe met de rij van elke \omterm{def:b3:point-groups:irrep}{irreducibele
        representatie}, en weeg elke klasse met haar grootte.
  \item Controleer: de resultaten zijn natuurlijke getallen, en $\sum_in_id_i$
        is gelijk aan het \omterm{def:b3:point-groups:representation}{karakter} onder $E$ (de grootte van de basis).
\end{enumerate}
\end{method}

\section{Symmetrieaangepaste orbitalen}

\begin{definition}[Projectieoperator]\label{def:b3:group-theory-applied:projection}
De \emph{projectieoperator}\index{projectieoperator} op de \omterm{def:b3:point-groups:irrep}{irreducibele
representatie} $\Gamma_i$ is
\[
\hat P_i = \frac{d_i}{h}\sum_R\chi_i(R)^*\,\hat R,
\]
waarin $\hat R$ de operatie $R$ op een functie toepast.
\end{definition}

\begin{proposition}[Wat de projectieoperator doet]\label{prop:b3:group-theory-applied:projection}
Toegepast op een willekeurige functie $f$ is $\hat P_if$ ofwel nul, ofwel
een functie die transformeert als $\Gamma_i$ (een symmetrieaangepaste
combinatie); voor een
eendimensionale $\Gamma_i$ blijft ze, op het teken $\chi_i(S)$ na,
onveranderd onder elke operatie $S$.
\end{proposition}

\begin{proof}[Gedeeltelijk bewijs]
Voor $d_i = 1$: $\hat S\hat P_if = \frac1h\sum_R\chi_i(R)\hat S\hat Rf$. Stel $R' = SR$,
dat de hele groep doorloopt wanneer $R$ dat doet; $\chi_i(R) = \chi_i(S^{-1}R') =
\chi_i(S)\chi_i(R')$ voor een eendimensionale representatie (de \omterm{def:b3:point-groups:representation}{karakters}
zijn de matrices zelf). Dus $\hat S\hat P_if = \chi_i(S)\hat P_if$. Het
algemene geval gebruikt de grote orthogonaliteitsstelling en wordt hier
aangenomen.
\end{proof}

\begin{method}[Symmetrieaangepaste combinaties opbouwen]\label{met:b3:group-theory-applied:salc}
\begin{enumerate}
  \item Kies één functie van de verzameling, $h_1$, en zet in een tabel naar
        welke functie elke operatie haar stuurt.
  \item Vermenigvuldig voor elke \omterm{def:b3:point-groups:irrep}{irreducibele representatie} elk beeld met
        het \omterm{def:b3:point-groups:representation}{karakter} van de operatie, tel op en normeer (met verwaarlozing
        van de overlap tussen de atoomfuncties).
  \item Projecteer voor een ontaarde representatie een tweede functie om
        een tweede lid te verkrijgen, en orthogonaliseer daarna.
\end{enumerate}
\end{method}

\begin{example}[De combinaties van ammoniak en water]\label{ex:b3:group-theory-applied:salcs}
In \ce{NH3} is $\hat P_{A_1}h_1 \propto h_1 + h_2 + h_3$ (alle \omterm{def:b3:point-groups:representation}{karakters} 1). Voor $E$
(\omterm{def:b3:point-groups:representation}{karakters} $2, -1, 0$) is $\hat P_Eh_1 \propto 2h_1 - h_2 - h_3$; $h_2$
projecteren en
orthogonaliseren geeft de partner $h_2 - h_3$. Genormeerd: $a_1 =
\frac{1}{\sqrt3}(h_1 + h_2 + h_3)$, $e = \frac{1}{\sqrt6}(2h_1 - h_2 - h_3)$,
$\frac{1}{\sqrt2}(h_2 - h_3)$. In \ce{H2O}, geplaatst in het $xz$-vlak, is $h_1 + h_2$
$a_1$ en $h_1 - h_2$ $b_1$ (het verandert van teken onder $C_2$ en onder $\sigma_v'(yz)$,
die de waterstofatomen verwisselen). Met het molecuul in het $yz$-vlak,
zoals sommige boeken verkiezen, zijn de labels $b_1$ en $b_2$ overal
verwisseld.
\end{example}

\begin{omfigure}
\begin{tikzpicture}[font=\small]
  \foreach \x/\lab/\ca/\cb/\cc in {0/{$a_1$}/1/1/1, 4/{$e$}/2/-1/-1, 8/{$e$}/0/1/-1}{
    \begin{scope}[xshift=\x cm]
      \foreach \a/\c in {90/\ca,210/\cb,330/\cc}{
        \pgfmathsetmacro{\r}{0.18+0.13*abs(\c)}
        \ifdim \c pt > 0pt \fill[omDef!55, draw=omDef] (\a:1.0) circle (\r);
        \else \ifdim \c pt < 0pt \fill[omThm!25, draw=omThm] (\a:1.0) circle (\r);
        \else \draw[gray, densely dotted] (\a:1.0) circle (0.12); \fi\fi
      }
      \node[aN, minimum size=4.6mm] at (0,0) {};
      \node[anchor=north] at (0,-1.45) {\lab};
    \end{scope}}
  \node[anchor=north, font=\footnotesize] at (0,-1.85) {$h_1 + h_2 + h_3$};
  \node[anchor=north, font=\footnotesize] at (4,-1.85) {$2h_1 - h_2 - h_3$};
  \node[anchor=north, font=\footnotesize] at (8,-1.85) {$h_2 - h_3$};
\end{tikzpicture}

{\small De symmetrieaangepaste combinaties van de drie $1s$-orbitalen van
waterstof in ammoniak, gezien langs de $C_3$-as ($h_1$ bovenaan). Blauw:
positieve coëfficiënt, rood: negatieve; de grootte van elke schijf groeit
met de coëfficiënt; gestippeld: nul.}
\end{omfigure}

\begin{example}[Zes liganden rond een metaal]\label{ex:b3:group-theory-applied:ml6}
Voor zes $\sigma$-orbitalen van liganden die vanuit de hoekpunten van een
octaëder naar een metaal wijzen, zijn de \omterm{def:b3:point-groups:representation}{karakters} op de klassen van $O_h$ ($E$, $8C_3$, $6C_2$, $6C_4$,
$3C_2$, $i$, $6S_4$, $8S_6$, $3\sigma_h$, $6\sigma_d$) gelijk aan $6, 0, 0, 2, 2, 0, 0, 0, 4,
2$, en de reductie geeft $\Gamma_\sigma = A_{1g} + E_g + T_{1u}$. De orbitalen
$s$
($a_{1g}$), $d_{z^2}$ en $d_{x^2-y^2}$ ($e_g$) en $p$ ($t_{1u}$) van het metaal
vinden partners bij de liganden; zijn $d_{xy}$, $d_{xz}$, $d_{yz}$
($t_{2g}$) vinden er geen en blijven niet-bindend
voor $\sigma$ --- de oorsprong van de opsplitsing $t_{2g}$/$e_g$ uit het
deel van jaar~2,
uitgewerkt in \cref{ch:b3:organometallic-bonding,ch:b3:complex-spectra-magnetism}.
\end{example}

\begin{omfigure}
\begin{tikzpicture}[font=\small, yscale=0.95]
  % O orbitals (left), MOs (centre), H combinations (right); schematic energies
  \foreach \y/\l in {0/{$2s\ (a_1)$}}{\draw[omstate] (0,\y) -- (1.2,\y); \node[anchor=east] at (-0.05,\y) {\l};}
  \draw[omstate] (0,3.0) -- (0.36,3.0) (0.42,3.0) -- (0.78,3.0) (0.84,3.0) -- (1.2,3.0);
  \node[anchor=east] at (-0.05,3.0) {$2p\ (b_1, a_1, b_2)$};
  \draw[omstate] (6.8,3.9) -- (8.0,3.9); \node[anchor=west] at (8.05,3.9) {$h_1 - h_2\ (b_1)$};
  \draw[omstate] (6.8,3.6) -- (8.0,3.6); \node[anchor=west] at (8.05,3.5) {$h_1 + h_2\ (a_1)$};
  \foreach \y/\l in {-0.6/{$2a_1$}, 1.6/{$1b_1$}, 2.4/{$3a_1$}, 3.0/{$1b_2$}, 5.6/{$4a_1$}, 6.2/{$2b_1$}}{
    \draw[omstate] (3.4,\y) -- (4.6,\y); \node[anchor=west] at (4.65,\y) {\l};}
  \foreach \y in {-0.6,1.6,2.4,3.0}{\node[font=\scriptsize] at (4.0,\y+0.16) {$\uparrow\downarrow$};}
  \draw[densely dotted, gray] (1.2,0) -- (3.4,-0.6) (1.2,3.0) -- (3.4,1.6) (1.2,3.0) -- (3.4,2.4) (1.2,3.0) -- (3.4,3.0)
     (6.8,3.6) -- (4.6,-0.6) (6.8,3.9) -- (4.6,1.6) (6.8,3.6) -- (4.6,2.4) (1.2,3.0) -- (3.4,6.2) (6.8,3.9) -- (4.6,6.2) (6.8,3.6) -- (4.6,5.6);
  \node[anchor=north] at (0.6,-1.0) {O};
  \node[anchor=north] at (4.0,-1.0) {\ce{H2O}};
  \node[anchor=north] at (7.4,-1.0) {2 H};
  \draw[->] (-2.3,-0.9) -- (-2.3,6.4) node[above] {$E$};
\end{tikzpicture}

{\small Moleculaire valentieorbitalen van water (schematische energieën,
molecuul in het $xz$-vlak). Alleen orbitalen met dezelfde symmetrie
mengen: de $b_1$-combinatie van waterstof met $2p_x$, de $a_1$-combinatie
met $2s$ en $2p_z$; $2p_y$ ($b_2$) heeft geen partner en blijft een
niet-bindend vrij elektronenpaar, de hoogste bezette orbitaal.}
\end{omfigure}

\section{Normaaltrillingen}

\begin{definition}[Normaaltrilling]\label{def:b3:group-theory-applied:normal-mode}
Een \emph{normaaltrilling}\index{normaaltrilling} van een molecuul is een
collectieve vibratie waarin alle atomen met dezelfde frequentie en in fase
oscilleren, langs de eigenvector van de massagewogen hessematrix
(\cref{prop:b3:computational-chemistry:hessian}). Elke normaaltrilling
transformeert als een \omterm{def:b3:point-groups:irrep}{irreducibele representatie} van de \omterm{def:b3:point-groups:point-group}{puntgroep}.
\end{definition}

\begin{proposition}[De representatie van alle bewegingen]\label{prop:b3:group-theory-applied:gamma-3n}
Hecht aan elk atoom drie verplaatsingsvectoren $x$, $y$, $z$. Onder een
operatie $R$ is het \omterm{def:b3:point-groups:representation}{karakter} van deze $3N$-dimensionale representatie
\[
\chi_{3N}(R) = (\text{aantal atomen dat op zijn plaats blijft}) \times \chi_{xyz}(R),
\]
met $\chi_{xyz} = 1 + 2\cos\theta$ voor een rotatie over $\theta$ en $-1 + 2\cos\theta$
voor een oneigenlijke rotatie (dus 3 voor $E$, $-1$ voor $C_2$, 0 voor $C_3$, 1 voor $\sigma$,
$-3$ voor $i$).
\end{proposition}

\begin{proof}
Een atoom dat naar een andere positie verplaatst wordt, neemt zijn
vectoren mee buiten de diagonaal van de matrix: het draagt niets bij aan
het spoor. Bij een atoom dat op zijn plaats blijft, worden de drie
vectoren getransformeerd door de $3\times3$-matrix van $R$, waarvan het
spoor berekend wordt in een assenstelsel met $z$ langs de as: $\cos\theta + \cos\theta + 1$ voor een
rotatie, waarbij het laatste element $-1$ wordt voor een oneigenlijke.
\end{proof}

\begin{proposition}[Vibraties tellen]\label{prop:b3:group-theory-applied:mode-count}
$\Gamma_{\mathrm{vib}} = \Gamma_{3N} - \Gamma_{\mathrm{trans}} - \Gamma_{\mathrm{rot}}$, waarin
$\Gamma_{\mathrm{trans}}$ de representatie van $(x, y, z)$ is en $\Gamma_{\mathrm{rot}}$ die
van $(R_x, R_y, R_z)$; een niet-lineair molecuul heeft $3N - 6$ vibraties,
een lineair
$3N - 5$.
\end{proposition}

\begin{proof}
De $3N$ verplaatsingen beschrijven elke beweging: drie translaties van het
massamiddelpunt, drie rotaties (twee voor een lineair molecuul, omdat een
draaiing om de eigen as geen enkele kern verplaatst), en de vibraties.
Deze deelruimten mengen niet onder de operaties, dus hun representaties
tellen op.
\end{proof}

\begin{example}[Water]\label{ex:b3:group-theory-applied:water}
% ledger: vib:H2O.sym, vib:H2O.bend, vib:H2O.asym
\ce{H2O} in $C_{2v}$, molecuul in het $xz$-vlak. Atomen die op hun plaats
blijven: 3, 1, 3, 1;
$\chi_{xyz}$: $3, -1, 1, 1$; dus $\chi_{3N} = 9, -1, 3, 1$. De reductie geeft $3A_1 +
A_2 + 3B_1 + 2B_2$. Na aftrek van de translaties $A_1 + B_1 + B_2$ en de
rotaties $A_2 + B_1 + B_2$ blijft
over: $\Gamma_{\mathrm{vib}} = 2A_1 + B_1$. De twee $a_1$-trillingen zijn de
symmetrische strekvibratie (\qty{3657}{cm^{-1}}) en de buigvibratie
(\qty{1595}{cm^{-1}}); de $b_1$-trilling is de antisymmetrische
strekvibratie (\qty{3756}{cm^{-1}}).
\end{example}

\begin{omfigure}
\begin{tikzpicture}[font=\small, >=Stealth]
  \foreach \x/\name/\lab in {0/sym/{$\nu_1$, $a_1$,\\ symmetrische strekvibratie}, 4.3/bend/{$\nu_2$, $a_1$,\\ buigvibratie}, 8.6/asym/{$\nu_3$, $b_1$,\\ antisymmetrische strekvibratie}}{
    \begin{scope}[xshift=\x cm]
      \node[aO] (o\name) at (0,0.35) {};
      \node[aH] (a\name) at (-0.8,-0.25) {};
      \node[aH] (b\name) at (0.8,-0.25) {};
      \begin{scope}[on background layer]\draw[bond] (o\name.center) -- (a\name.center) (o\name.center) -- (b\name.center);\end{scope}
      \node[anchor=north, align=center, font=\footnotesize] at (0,-0.9) {\lab};
    \end{scope}}
  % symmetric stretch: H outward along the bonds, O up a little
  \draw[->, omThm, thick] (-0.95,-0.36) -- (-1.42,-0.71);
  \draw[->, omThm, thick] (0.95,-0.36) -- (1.42,-0.71);
  \draw[->, omThm, thick] (0,0.74) -- (0,0.98);
  % bend: each H moves perpendicular to its bond, closing the angle; O up
  \draw[->, omThm, thick] (4.3-0.92,-0.36) -- (4.3-0.62,-0.76);
  \draw[->, omThm, thick] (4.3+0.92,-0.36) -- (4.3+0.62,-0.76);
  \draw[->, omThm, thick] (4.3,0.74) -- (4.3,0.98);
  % antisymmetric stretch: one H out, the other in, O sideways
  \draw[->, omThm, thick] (8.6-0.95,-0.36) -- (8.6-1.42,-0.71);
  \draw[->, omThm, thick] (8.6+1.02,0.0) -- (8.6+0.70,0.24);
  \draw[->, omThm, thick] (8.6-0.05,0.80) -- (8.6+0.25,0.80);
\end{tikzpicture}

{\small De drie \omterm{def:b3:group-theory-applied:normal-mode}{normaaltrillingen} van water (pijlen: verplaatsingen, niet op
schaal). Beide $a_1$-trillingen behouden de volledige symmetrie van het
molecuul; de $b_1$-trilling verandert van teken onder $C_2$.}
\end{omfigure}

\begin{example}[Nog vier moleculen]\label{ex:b3:group-theory-applied:table}
Dezelfde procedure (berekend en gecontroleerd in de figuurgegevens van dit
hoofdstuk) geeft:
\begin{center}\small
\begin{tabular}{llll}
\hline
molecuul & groep & $\Gamma_{\mathrm{vib}}$ & trillingen \\ \hline
\ce{NH3} & $C_{3v}$ & $2A_1 + 2E$ & 6 \\
\ce{CH4} & $T_d$ & $A_1 + E + 2T_2$ & 9 \\
\ce{CO2} & $D_{\infty h}$ (via $D_{2h}$) & $\Sigma_g^+ + \Sigma_u^+ + \Pi_u$ ($A_g + B_{1u} + B_{2u} + B_{3u}$) & 4 \\
\ce{BF3} & $D_{3h}$ & $A_1' + 2E' + A_2''$ & 6 \\
\ce{XeF4} & $D_{4h}$ & $A_{1g} + B_{1g} + B_{2g} + A_{2u} + B_{2u} + 2E_u$ & 9 \\
\ce{SF6} & $O_h$ & $A_{1g} + E_g + T_{2g} + 2T_{1u} + T_{2u}$ & 15 \\ \hline
\end{tabular}
\end{center}
Voor een lineair molecuul wordt de oneindige groep vervangen door haar
\omterm{def:b3:point-groups:class}{ondergroep} $D_{2h}$,
waarin de ontaarde buigvibratie $\Pi_u$ opsplitst in $B_{2u} + B_{3u}$.
\end{example}

\section{Selectieregels}

\begin{definition}[Direct product van representaties]\label{def:b3:group-theory-applied:direct-product}
Het \emph{direct product}\index{direct product} $\Gamma_i\otimes\Gamma_j$ van twee
representaties is de representatie die gedragen wordt door de producten
van hun basisfuncties; haar \omterm{def:b3:point-groups:representation}{karakters} zijn de producten $\chi_i(R)\chi_j(R)$.
\end{definition}

\begin{theorem}[Integralen die nul zijn]\label{thm:b3:group-theory-applied:vanishing-integral}
Een integraal $\int f_1f_2f_3\,\dd\tau$ over de hele ruimte kan alleen
verschillen van nul als het
\omterm{def:b3:group-theory-applied:direct-product}{direct product} $\Gamma_1\otimes\Gamma_2\otimes\Gamma_3$ de \omterm{def:b3:group-theory-applied:reducible}{totaalsymmetrische
representatie} bevat.
\end{theorem}

\begin{proof}
Een \omterm{def:b3:point-groups:operation}{symmetrieoperatie} hernummert alleen de punten van de ruimte, zodat de
waarde van de integraal onveranderd blijft wanneer de integrand door een
willekeurige $R$ getransformeerd wordt. Neem het gemiddelde van de
getransformeerde integrand over de groep: het gemiddelde van de componenten van
de integrand die bij elke \omterm{def:b3:point-groups:irrep}{irreducibele representatie} horen, is
$\frac1h\sum_R\hat R(\cdot)$, de projectie op de \omterm{def:b3:group-theory-applied:reducible}{totaalsymmetrische
representatie} (\cref{def:b3:group-theory-applied:projection} met alle
$\chi = 1$). Elke andere component heeft gemiddelde nul; als het product
geen totaalsymmetrisch deel bevat, is de integraal nul.
\end{proof}

\begin{definition}[IR-actief, Raman-actief]\label{def:b3:group-theory-applied:activity}
Een fundamentele vibratie is \emph{IR-actief}\index{IR-actief} als ze
infraroodlicht kan absorberen, en \emph{Raman-actief}\index{Raman-actief}
als ze in het
ramanspectrum verschijnt (\cref{ch:b3:rovibrational-spectroscopy}).
\end{definition}

\begin{corollary}[Selectieregels voor infrarood en Raman]\label{cor:b3:group-theory-applied:ir-raman}
Een fundamentele overgang (van het grondniveau, totaalsymmetrisch, naar
één kwantum van een trilling met symmetrie $\Gamma$) is alleen \omterm{def:b3:group-theory-applied:activity}{IR-actief}
als $\Gamma$ de
representatie van $x$, $y$ of $z$ is, en alleen \omterm{def:b3:group-theory-applied:activity}{Raman-actief} als $\Gamma$ de
representatie van een kwadratische functie ($x^2$, $xy$, \dots) is.
\end{corollary}

\begin{proof}
De absorptie-intensiteit hangt af van $\int\psi_0\,\mu_k\,\psi_1\dd\tau$ met de
dipoolcomponenten $\mu_k$, die transformeren als $x$, $y$, $z$; $\psi_0$ is
totaalsymmetrisch en $\psi_1$ transformeert als $\Gamma$. Volgens de
stelling kan de integraal alleen verschillen van nul
als $\Gamma\otimes\Gamma_{\mu_k}$ $A_1$ bevat, en dat gebeurt alleen als $\Gamma =
\Gamma_{\mu_k}$ (het product van twee \omterm{def:b3:point-groups:irrep}{irreducibele representaties} bevat de
totaalsymmetrische alleen als ze gelijk zijn, voor reële \omterm{def:b3:point-groups:representation}{karakters}).
Ramanverstrooiing hangt af van de componenten $\alpha_{kl}$ van de
polariseerbaarheid, die
transformeren als de producten $kl$.
\end{proof}

\begin{proposition}[Regel van wederzijdse uitsluiting]\label{prop:b3:group-theory-applied:mutual-exclusion}
In een molecuul met een \omterm{def:b3:point-groups:inversion}{inversiecentrum} is geen enkele fundamentele
overgang tegelijk IR- en \omterm{def:b3:group-theory-applied:activity}{Raman-actief}: de \emph{regel van wederzijdse
uitsluiting}\index{regel van wederzijdse uitsluiting}.
\end{proposition}

\begin{proof}
Met een \omterm{def:b3:point-groups:inversion}{inversiecentrum} is elke \omterm{def:b3:point-groups:irrep}{irreducibele representatie} $g$ of $u$. De
coördinaten $x$, $y$, $z$ veranderen van teken onder $i$ (ze zijn $u$);
kwadratische functies niet ($g$). Een trilling is alleen \omterm{def:b3:group-theory-applied:activity}{IR-actief} als ze
$u$ is, en alleen \omterm{def:b3:group-theory-applied:activity}{Raman-actief} als ze
$g$ is.
\end{proof}

\begin{example}[Koolstofdioxide en het broeikaseffect]\label{ex:b3:group-theory-applied:co2}
% ledger: vib:CO2.sym, vib:CO2.bend, vib:CO2.asym, ir:CO2.asym.int, ir:CO2.bend.int
\ce{CO2} heeft een \omterm{def:b3:point-groups:inversion}{inversiecentrum}. Zijn symmetrische strekvibratie ($\Sigma_g^+$,
\qty{1333}{cm^{-1}}) is alleen \omterm{def:b3:group-theory-applied:activity}{Raman-actief}; de antisymmetrische
strekvibratie ($\Sigma_u^+$,
\qty{2349}{cm^{-1}}) en de buigvibratie ($\Pi_u$, \qty{667}{cm^{-1}}) zijn
alleen \omterm{def:b3:group-theory-applied:activity}{IR-actief}.
De buigvibratie absorbeert rond \qty{15}{\micro m}, dicht bij het maximum
van de thermische straling van de aarde: daarom is koolstofdioxide een
broeikasgas. \ce{N2} en
\ce{O2} hebben één enkele vibratie, $\Sigma_g^+$: IR-inactief.
\end{example}

\begin{omfigure}
\begin{tikzpicture}
\begin{axis}[omir, width=0.92\linewidth, height=4.4cm, xmin=400, xmax=2600,
  ymin=0, ymax=1.3, ytick=\empty, xlabel={golfgetal / cm$^{-1}$},
  ylabel={signaal}, legend cell align=left,
  legend style={font=\small, at={(0.98,0.95)}, anchor=north east, draw=none}]
\addplot[omDef, thick] table[x=nu, y=ir] {figdata/out/bachelor-3-co2-spectra.dat};
\addlegendentry{infraroodabsorptie}
\addplot[omThm, thick, dashed] table[x=nu, y=raman] {figdata/out/bachelor-3-co2-spectra.dat};
\addlegendentry{ramanverstrooiing}
\node[font=\small, anchor=south] at (axis cs:2349,1.02) {$\Sigma_u^+$};
\node[font=\small, anchor=south] at (axis cs:667,0.1) {$\Pi_u$};
\node[font=\small, anchor=south, omThm] at (axis cs:1333,1.02) {$\Sigma_g^+$};
\end{axis}
\end{tikzpicture}

{\small Synthetische spectra van gasvormig \ce{CO2} (bandposities en
relatieve infraroodintensiteiten uit de gemeten waarden; vormen
schematisch, zonder rotatiestructuur). Wederzijdse uitsluiting: geen enkele
band komt in beide voor.}
\end{omfigure}

\begin{method}[Een vibratiespectrum voorspellen]\label{met:b3:group-theory-applied:vibrations}
\begin{enumerate}
  \item Zoek de \omterm{def:b3:point-groups:point-group}{puntgroep}; bereken $\chi_{3N}$ uit de atomen die op hun
        plaats blijven.
  \item Reduceer, trek translaties en rotaties af: $\Gamma_{\mathrm{vib}}$.
  \item Markeer met de kolommen rechts in de tabel elke trilling als
        \omterm{def:b3:group-theory-applied:activity}{IR-actief} ($x$, $y$, $z$) en/of \omterm{def:b3:group-theory-applied:activity}{Raman-actief} (kwadratisch); een
        trilling die geen van beide is, is stil.
  \item Tel de banden: één per actieve trilling, waarbij een ontaarde
        trilling één band geeft.
\end{enumerate}
\end{method}

\begin{method}[CO-strekvibraties tellen]\label{met:b3:group-theory-applied:co-count}
Neem in een metaalcarbonyl de $n$ bindingsvectoren C--O als basis (een
atoom dat op zijn plaats blijft telt 1, al de rest niets): reduceer om
$\Gamma_{\mathrm{CO}}$ te verkrijgen, en pas daarna
de regels voor IR en Raman toe. Het aantal sterke banden rond
\qtyrange{1850}{2150}{cm^{-1}} identificeert het isomeer.
\end{method}

\begin{example}[Cis en trans]\label{ex:b3:group-theory-applied:cis-trans}
cis-\ce{ML4(CO)2}, $C_{2v}$: \omterm{def:b3:point-groups:representation}{karakters} $2, 0, 2, 0$, $\Gamma_{\mathrm{CO}} = A_1 + B_1$,
twee IR-banden. trans-\ce{ML4(CO)2}, $D_{4h}$: $\Gamma_{\mathrm{CO}} = A_{1g} + A_{2u}$, één
IR-band ($A_{2u}$) en één ramanband ($A_{1g}$). Eén enkele CO-band in het
infrarood betekent trans.
\end{example}

\section{Toepassingen op elektronische overgangen}

Dezelfde stelling geldt voor elektronische overgangen: een elektronische
overgang tussen de toestanden $\Psi_1$ en $\Psi_2$ is alleen toegestaan
als $\Gamma_1\otimes\Gamma_2$
de representatie van $x$, $y$ of $z$ bevat, en het licht ervan is
gepolariseerd langs die as.

\begin{example}[Overgangen van water en formaldehyde]\label{ex:b3:group-theory-applied:electronic}
In water ($C_{2v}$, $xz$-vlak) geeft de promotie van een elektron van $1b_2$ (het vrije elektronenpaar)
naar $4a_1$ een aangeslagen toestand met symmetrie $B_2\otimes A_1 = B_2$,
die $y$ is: toegestaan, gepolariseerd loodrecht op het molecuulvlak. In
formaldehyde
($C_{2v}$) gaat de overgang $n\to\pi^*$ van een vrij elektronenpaar van
zuurstof in het vlak ($b_1$)
naar de $\pi^*$-orbitaal, opgebouwd uit $p$-orbitalen loodrecht op het vlak
($b_2$; molecuul in het $xz$-vlak, C=O langs $z$): $B_1\otimes B_2 = A_2$,
dat geen van $x$, $y$, $z$ is: symmetrieverboden, en daarom is de band
ervan (\cref{ch:b3:electronic-spectroscopy}) zo zwak.
\end{example}

\begin{history}[Wigner, Bethe en de symmetrie van moleculen]
De groepentheorie deed haar intrede in de kwantummechanica met het werk
van Eugene Wigner over atoomspectra (1927--31) en het artikel van Hans
Bethe over de opsplitsing van atoomtermen in kristallen (1929). E. Bright
Wilson paste haar in 1934 toe op moleculaire vibraties; Robert Mulliken
gaf de labels die nog steeds voor orbitalen en toestanden gebruikt worden.
\end{history}

\begin{inthelab}[Infrarood en Raman naast elkaar]
Een infraroodspectrometer meet het licht dat een monster absorbeert; een
ramanspectrometer belicht het met een laser en analyseert het zwakke
verstrooide licht bij verschoven golflengten. Beide spectra van dezelfde
verbinding opnemen is de klassieke test op een symmetriecentrum: als geen
enkele band samenvalt, is het molecuul waarschijnlijk centrosymmetrisch.
Ramanlasers zijn van klasse 3B of 4: de bundel wordt afgeschermd en de
gebruikers dragen een veiligheidsbril die geschikt is voor de golflengte.
\end{inthelab}

\section{Oefeningen}

\begin{exercise}[$\star$]\label{exo:b3:group-theory-applied:1}
Reduceer de representatie van $C_{2v}$ met \omterm{def:b3:point-groups:representation}{karakters} $(5, -1, 3, 1)$ op $(E, C_2,
\sigma_v(xz), \sigma_v'(yz))$. Controleer de dimensie.
\end{exercise}

\begin{exercise}[$\star$]\label{exo:b3:group-theory-applied:2}
\ce{SO2} is gebogen, zoals water. Geef $\Gamma_{\mathrm{vib}}$ en zeg welke
trillingen IR- en \omterm{def:b3:group-theory-applied:activity}{Raman-actief} zijn.
\end{exercise}

\begin{exercise}[$\star$]\label{exo:b3:group-theory-applied:3}
Bereken in $C_{2v}$ de directe producten $B_1\otimes B_2$, $A_2\otimes B_1$, en
$E\otimes E$ in $C_{3v}$ (reduceer dat laatste).
\end{exercise}

\begin{exercise}[$\star$]\label{exo:b3:group-theory-applied:4}
% ledger: vib:CH4.nu1, vib:CH4.nu2, vib:CH4.nu3, vib:CH4.nu4
\ce{CH4} heeft fundamentele banden bij 2917 ($a_1$), 1534 ($e$), 3019 en 1306 (beide $t_2$)
\unit{cm^{-1}}. Welke verschijnen in het infraroodspectrum, welke in het
ramanspectrum?
\end{exercise}

\begin{exercise}[$\star\star$]\label{exo:b3:group-theory-applied:5}
Leid voor \ce{BF3} $\Gamma_{\mathrm{vib}} = A_1' + 2E' + A_2''$ af uit de atomen
die op hun plaats blijven, en tel daarna de IR- en ramanbanden. Welke
trilling ziet men in slechts één van de twee spectra, en welke in beide?
\end{exercise}

\begin{exercise}[$\star\star$]\label{exo:b3:group-theory-applied:6}
Hoeveel IR- en ramanbanden vertoont \ce{SF6}? Welke trilling is stil? Waarom
verschijnt geen enkele band in beide spectra?
\end{exercise}

\begin{exercise}[$\star\star$]\label{exo:b3:group-theory-applied:7}
Pas de \omterm{def:b3:group-theory-applied:projection}{projectieoperator} toe op $h_2$ voor de representatie $E$ van $C_{3v}$ en
toon aan dat het resultaat, orthogonaal gemaakt ten opzichte van $2h_1 - h_2 - h_3$, $h_2 - h_3$ geeft.
\end{exercise}

\begin{exercise}[$\star\star$]\label{exo:b3:group-theory-applied:8}
Voorspel het aantal CO-banden in het IR-spectrum van fac- en
mer-\ce{ML3(CO)3}, en van
\ce{M(CO)6}.
\end{exercise}

\begin{exercise}[$\star\star$]\label{exo:b3:group-theory-applied:9}
Toon aan dat de zes $\sigma$-orbitalen van de liganden van een octaëdrisch
complex
$A_{1g} + E_g + T_{1u}$ overspannen, en zeg met welke metaalorbitalen elk
ervan kan combineren.
\end{exercise}

\begin{exercise}[$\star\star\star$]\label{exo:b3:group-theory-applied:10}
De reguliere representatie heeft $\chi(E) = h$ en $\chi(R) = 0$ voor $R \ne E$. Toon
met de \omterm{thm:b3:group-theory-applied:reduction}{reductieformule} aan dat ze elke \omterm{def:b3:point-groups:irrep}{irreducibele
representatie} $\Gamma_i$ precies $d_i$ keer bevat, en leid $\sum_id_i^2 = h$ af.
\end{exercise}

\begin{exercise}[$\star\star\star$]\label{exo:b3:group-theory-applied:11}
Acetyleen \ce{HC#CH} is lineair. Werk in $D_{2h}$, bepaal $\Gamma_{\mathrm{vib}}$,
hergroepeer het
in labels van $D_{\infty h}$, en voorspel de IR- en ramanspectra.
\end{exercise}

\begin{exercise}[$\star\star\star$]\label{exo:b3:group-theory-applied:12}
Welke van deze promoties van één elektron in water ($xz$-vlak) geven
toegestane overgangen, en met welke polarisatie: $1b_2 \to 4a_1$, $3a_1 \to 4a_1$,
$1b_2 \to 2b_1$, $1b_1 \to 4a_1$?
\end{exercise}

\section{Probleem: cis of trans? Een carbonylcomplex gelezen aan zijn spectrum}

\begin{problem}[{Weekendprobleem --- puntgroepen en representaties van de
CO-strekvibraties van vier isomeren, hun infrarood- en ramanbanden, en de
identificatie van een product aan de hand van zijn gemeten spectrum}]
\label{pb:b3:group-theory-applied:1}
% ledger: ct:C2v, ct:C3v, ct:D4h
Een octaëdrisch metaal M draagt carbonylliganden en andere liganden L (elk
L geteld als een punt). Vier isomeren worden bekeken: cis- en
trans-\ce{ML4(CO)2}, en fac- en mer-\ce{ML3(CO)3}. Neem aan dat de liganden L
de symmetrie niet verder verlagen dan hun posities opleggen. Een synthese
van
\ce{ML3(CO)3} levert een product waarvan het infraroodspectrum drie sterke
banden vertoont bij 2048, 1965 en \qty{1938}{cm^{-1}} (gegevens van het
probleem); zijn ramanspectrum vertoont drie banden op vrijwel dezelfde
posities.

\medskip\noindent\textbf{Deel I --- \omterm{def:b3:point-groups:point-group}{Puntgroepen}.}
\begin{enumerate}
  \item Teken de vier isomeren op een octaëder.
  \item Geef de \omterm{def:b3:point-groups:point-group}{puntgroep} van cis-\ce{ML4(CO)2}.
  \item Van trans-\ce{ML4(CO)2}.
  \item Van fac-\ce{ML3(CO)3}.
  \item Van mer-\ce{ML3(CO)3}.
  \item Welke van de vier hebben een \omterm{def:b3:point-groups:inversion}{inversiecentrum}?
\end{enumerate}

\medskip\noindent\textbf{Deel II --- De CO-strekvibraties.}
\begin{enumerate}[resume]
  \item Leg uit waarom de bindingsvectoren C--O een representatie dragen
        waarin een operatie 1 bijdraagt per CO dat op zijn plaats blijft.
  \item Bepaal $\Gamma_{\mathrm{CO}}$ van het cis-isomeer.
  \item Van het trans-isomeer.
  \item Van het fac-isomeer.
  \item Van het mer-isomeer.
  \item Controleer elke dimensie aan de hand van het aantal CO-liganden.
\end{enumerate}

\medskip\noindent\textbf{Deel III --- Activiteit.}
\begin{enumerate}[resume]
  \item Tel de IR-actieve CO-strekvibraties van elk isomeer.
  \item Tel de Raman-actieve.
  \item Welk isomeer volgt de \omterm{prop:b3:group-theory-applied:mutual-exclusion}{regel van wederzijdse uitsluiting}?
  \item Beschrijf in het trans-isomeer de IR-actieve trilling: strekken de
        twee CO in fase of in tegenfase?
  \item Waarom geven de drie CO in het fac-isomeer maar twee banden?
  \item Wat zou één enkele band (IR) je vertellen?
\end{enumerate}

\medskip\noindent\textbf{Deel IV --- Het product.}
\begin{enumerate}[resume]
  \item Op welk isomeer wijst het infraroodspectrum van het product?
  \item Is het ramanspectrum daarmee in overeenstemming?
  \item De hoogste band is de strekvibratie in fase van de CO-groepen. Tot
        welke \omterm{def:b3:point-groups:irrep}{irreducibele representatie} behoort ze in dat isomeer?
  \item Zou het spectrum van een mengsel van de twee isomeren kunnen komen?
        Welke bijkomende meting zou de knoop doorhakken?
  \item Formuleer het resultaat: het aantal IR-actieve CO-strekvibraties
        van het mer-isomeer, tegenover dat van het fac-isomeer.
\end{enumerate}
\end{problem}
